\section{Resultados e Discussão}
\label{sec:discussao}

\subsection{Critérios de apresentação e confronto das hipóteses}

Esta seção estabelece o roteiro de análise para a etapa posterior à qualificação. Ainda não foram estimados os modelos do artigo, e as pendências relativas às conversões brutas, à geografia histórica e às medidas institucionais impedem apresentar resultados como se o painel estivesse concluído.

A discussão começará pela descrição da amostra efetivamente utilizada: territórios, anos, vegetação inicial, conversões observadas, perdas por junção das bases e disponibilidade de cada indicador institucional. Em seguida, serão apresentados os coeficientes das exposições por produto e os efeitos marginais correspondentes, em pontos percentuais e hectares. Comparações entre capacidades produtivas baixas e altas utilizarão valores dentro do suporte dos dados, com intervalos de confiança e indicação do tamanho do choque de preço considerado.

A H1-a receberá apoio se os movimentos de preços se associarem a maior conversão nos territórios de maior capacidade para o uso beneficiado, sob as condições de identificação discutidas. Uma resposta nula poderá refletir ausência do mecanismo, baixa precisão ou inadequação da medida de exposição; essas possibilidades serão distinguidas pela amplitude dos intervalos, pela qualidade dos dados e pelos testes de sensibilidade. Um sinal contrário exigirá examinar substituição entre usos, intensificação e custos, sem alterar retrospectivamente a hipótese para acomodar o resultado.

Para H1-b, a discussão mostrará como a sensibilidade aos preços varia conforme o indicador institucional, sempre identificando o que ele mede. Um coeficiente positivo de autos de infração, por exemplo, não será interpretado como prova de que fiscalizar aumenta a conversão. Os autos podem acompanhar a ocorrência de infrações, e a relação estimada pode refletir essa seleção. A eventual identificação de uma mudança institucional será apresentada com seu grupo de comparação, vigência, pressupostos e limites de validade.

\subsection{Heterogeneidade e implicações econômicas}

A comparação entre biomas e atividades buscará identificar diferenças de resposta, e não apenas de níveis médios de conversão. Os resultados para municípios mistos serão mantidos visíveis. Quando o Pantanal não oferecer suporte para inferência separada, sua discussão será descritiva, sem extrapolação de coeficientes estimados na Amazônia ou no Cerrado.

O significado econômico dos efeitos será avaliado em relação à vegetação disponível e ao volume de conversão observado. Variações percentuais elevadas em territórios com conversão inicial muito pequena serão acompanhadas das áreas correspondentes. Também serão discutidos resultados sobre conversão para pastagem e agricultura e, quando observável, substituição entre áreas agropecuárias já estabelecidas. Uma expansão de produção em terra previamente convertida não será contabilizada como perda adicional de vegetação.

Em termos de política, diferenças territoriais poderão indicar que preços, custos de oportunidade e implementação institucional importam para a focalização dos recursos. Essa constatação não identificará, por si só, o arranjo ótimo de governança nem autorizará concluir que instrumentos nacionais ou globais são ineficazes. Tampouco permitirá comparar pagamentos por serviços ambientais, créditos de carbono e outros mecanismos sem dados sobre sua aplicação. O alcance da contribuição será discutir os condicionantes observados da conversão e as combinações institucionais compatíveis com a evidência.

\subsection{Limites de interpretação}

O painel municipal agrega agentes com condições distintas e não recupera diretamente expectativas individuais, custo de capital ou horizontes de decisão. Efeitos defasados indicarão persistência da resposta, mas não identificarão taxas de desconto. A série do IBAMA registra sanções administrativas federais; seus resultados não serão estendidos a decisões judiciais ou à totalidade da governança ambiental.

Por fim, a perda de vegetação aqui estudada não esgota a condição do capital natural: a conservação da cobertura pode coexistir com degradação de solo, fragmentação ou alteração hidrológica. O artigo identifica um elo da relação entre economia e meio ambiente, a formação do estoque pela decisão de uso da terra. A contribuição produtiva desse estoque será tratada no segundo artigo, preservando a autonomia dos dois testes.

